% --- Archon physics formalization source begin ---
% archon:physics
% archon:covers IPhO2026Problems/problem_ipho_2026_t1_b2.lean
% archon:source-report reports/ipho_2026/problem_ipho_2026_t1_b2.source.json
% archon:problem-id ipho_2026_t1
% archon:part-id T1-B2

\chapter{Ipho Physics problem ipho\_2026\_t1\_b2}
\label{ch:IPhO2026Problems_problem_ipho_2026_t1_b2}

\paragraph{Problem source.}
\#\# Physical scenario

At one instant a positron and an electron, each of mass m and charges of equal
magnitude and opposite sign, are separated by 100*a\_0.  Their velocities are
antiparallel and perpendicular to their separation.  Each particle has angular
momentum of magnitude mu*hbar about the center of mass.  The system is isolated,
classical, non-relativistic, and has only electrostatic interaction.  The Bohr
radius is a\_0 = 4*pi*epsilon\_0*hbar\textasciicircum{}2/(m*e\textasciicircum{}2), and k = 1/(4*pi*epsilon\_0).

\#\# Current subquestion T1-B2

For mu = 15/2 the pair is unbound. Find the angle between the asymptotic relative velocity u\_infinity and the initial positron line of motion.

\paragraph{Shared problem context.}
At one instant a positron and an electron, each of mass m and charges of equal
magnitude and opposite sign, are separated by 100*a\_0.  Their velocities are
antiparallel and perpendicular to their separation.  Each particle has angular
momentum of magnitude mu*hbar about the center of mass.  The system is isolated,
classical, non-relativistic, and has only electrostatic interaction.  The Bohr
radius is a\_0 = 4*pi*epsilon\_0*hbar\textasciicircum{}2/(m*e\textasciicircum{}2), and k = 1/(4*pi*epsilon\_0).

\paragraph{Current subquestion.}
For mu = 15/2 the pair is unbound. Find the angle between the asymptotic relative velocity u\_infinity and the initial positron line of motion.

\paragraph{Figure/image path.}
ipho\_2026\_source/image/T1\_page-2.png

\paragraph{Formalization target.}
create a compiling Lean file with sorry bodies at `IPhO2026Problems/problem\_ipho\_2026\_t1\_b2.lean`.
The Lean declarations must preserve the physical quantities, dimensions or dimensional roles, figure labels, governing-law hypotheses, and final relation expressed by this problem.
Use Mathlib/Physlib names found through LeanExplore where available. If a domain API is missing, introduce faithful local abstractions rather than scalar placeholder aliases.
This is an answer-blind solve-phase task. Derive a candidate result only from the problem-side evidence above; do not assume a target value, select a tolerance around a desired value, or encode a desired result into a candidate domain.
For numerical questions, define the raw end-to-end quantity and apply an explicit source-derived rounding rule. For identification questions, characterize or prove uniqueness before recording the derived candidate.

\begin{theorem}[Ipho Physics formalization target]
\label{thm:physics:ipho_2026_t1_b2:target}
The assigned autoformalize agent should translate this IPhO physics problem into Lean declarations in the covered file, with theorem and lemma proof bodies written as `by sorry`.
\end{theorem}
\begin{proof}
This is an autoformalization task, not a proof task. Produce faithful statements that can later be proved without weakening the source contract.
\end{proof}
% --- Archon physics formalization source end ---
